\documentclass[11pt]{article}
\usepackage[margin=1in]{geometry}
\usepackage{microtype}
\usepackage{amsmath,amssymb}
\usepackage[numbers,sort&compress]{natbib}
\usepackage{hyperref}
\usepackage{booktabs}
\usepackage{parskip}

\bibliographystyle{unsrtnat}

\title{Self/World Factorisation in Trained Agents:\\ Experimental Design}
\author{}
\date{}

\begin{document}
\maketitle

\noindent\emph{This is the working experimental-design document for the project.
The motivation, the developmental-cognitive-science case, and the broader
framing live in the grant version (\texttt{grant\_proposal.tex}); here
we keep only as much theory as is needed to fix what we build, what we measure,
and what each measurement would show. Two companion files state the geometric
and metacognitive targets the probes aim at: \texttt{strand\_two.tex} (the
geometry of the factorisation) and \texttt{strand\_three.tex} (epistemic actions
and introspective access).}

\section{Hypothesis, operationalised}

We study transformers trained on a turn-taking stream of actions and
observations generated by a coupled hidden process. The hypothesis has two
parts, each stated as a concrete measurable claim.

\begin{description}
  \item[H1 (belief).] The model represents a Bayes-adaptive posterior over the
  hidden \emph{World} state, linearly decodable from the residual stream and
  matching the analytic posterior --- including the action-conditioned update,
  in which the conditioning action is itself a sample from the model's own
  output distribution.
  \item[H2 (self).] The model represents a \emph{Self} component, separable from
  the belief, in one or more of three senses we keep distinct throughout:
  \emph{(a) token-type self} --- action vs.\ observation positions select
  different belief updates; \emph{(b) actor-identity self} --- a posterior over
  which policy produced the action stream, with the model's own policy a
  distinguished vertex; \emph{(c) epistemic self} --- a represented property of
  the model's \emph{own} belief state (e.g.\ its marginal uncertainties),
  causally consulted before the model acts to gain information.
\end{description}

The unifying geometric claim (H1$+$H2 together) is that the reachable-belief
manifold collapses onto a product $\Delta_S\times\Delta_W$ with independently
decodable marginals, rather than filling an entangled subset of the joint
simplex. This is made precise in \texttt{strand\_two.tex}.

\section{The SWITCH generative model}
\label{sec:switch}

We call the generative process the \textbf{SWITCH} (Self/World Interactive
Turn-taking Coupled HMM). It is small enough that the exact Bayes-adaptive
posterior is computable at every step, giving a ground-truth comparator the
language-model experiments lack.

\paragraph{State, action, observation spaces.} World hidden state
$w\in\mathcal{W}$ ($|\mathcal{W}|=n_W$); policy/self state $s\in\mathcal{S}$;
action $a\in\mathcal{A}$; observation $o\in\mathcal{O}$. All four are finite and
small (initial targets $n_W\in\{3,4,5\}$, $|\mathcal{A}|,|\mathcal{O}|\le 4$).

\paragraph{Dynamics.} The generative process is two coupled chains exchanging
visible tokens:
\begin{itemize}
  \item \textbf{World} chain: action-conditioned transition
  $w_{t+1}\sim T(\cdot\mid w_t,a_t)$ and emission $o_{t+1}\sim E(\cdot\mid
  w_{t+1})$ produced on the transition. (This half is an input--output HMM
  \citep{bengio1995iohmm} with actions as inputs.)
  \item \textbf{Self} chain: action emission $a_t\sim\pi(\cdot\mid s_t)$ and
  self-state update $s_{t+1}=g(s_t,a_t,o_{t+1})$.
\end{itemize}
The visible token stream is $x=(a_1,o_1,a_2,o_2,\dots,a_n,o_n)$: odd positions
carry Self-emitted actions, even positions carry World-emitted observations.
The actions and observations are the only coupling --- a Markov blanket between
Self and World. The joint $(T,E,\pi,g)$ is sampled once per SWITCH instance
(see \S\ref{sec:metalearn} for the meta-learning variant where it is resampled).
Optionally a scalar reward $r_t=R(w_{t+1},a_t)$ is defined, used only by the RL
regime.

\paragraph{Analytic ground truth (the comparator).} The exact Bayes-adaptive
World posterior $b_W^t(w)=P(w_t\mid x_{<t})$ is maintained by forward filtering
with two update maps selected by token parity:
\begin{align}
  \text{action } a_t:\quad
    & b_W(w)\;\propto\; b_W(w)\,\pi^{(c)}(a_t\mid w),
      &&\text{(actor-gated likelihood)} \label{eq:actorgate}\\
    & b_W \;\mapsto\; T(\cdot\mid\cdot,a_t)^{\top} b_W,
      &&\text{(control pushforward)} \label{eq:push}\\
  \text{observation } o_{t+1}:\quad
    & b_W(w')\;\propto\; E(o_{t+1}\mid w')\,b_W(w').
      &&\text{(filtering update)} \label{eq:filter}
\end{align}
The actor-gated factor \eqref{eq:actorgate} is what makes the active case
non-trivial: an action is evidence about $w$ only insofar as the acting policy
$\pi^{(c)}$ is state-dependent. For a state-\emph{independent} actor the factor
is constant and only the pushforward \eqref{eq:push} applies. Where multiple
actors are possible (Regime~5) the model must also maintain a posterior over the
actor identity $c$; where the World factors (Regime~6) the quotient/epimorphism
machinery of \texttt{strand\_three.tex} applies. The set of reachable $b_W^t$ is
the mixed-state presentation (MSP) we compare the decoded geometry against,
following \citep{shai2024belief}.

\paragraph{What we implement.} (i) A SWITCH sampler producing token streams and,
alongside each token, the analytic $b_W^t$, the analytic actor posterior (when
defined), and marginal entropies. (ii) A standard decoder-only transformer
trained on these streams. (iii) A probe/intervention harness operating on
residual-stream activations $h_\ell^t$. The same three components serve every
regime; regimes differ only in how the action stream is produced and where the
loss is placed.

\section{Training regimes}
\label{sec:regimes}

Regimes vary along two axes the hypothesis identifies --- whether the model
produces its own actions (passive vs.\ active), and whether the actor is the
model alone or one of a family (single- vs.\ multi-actor) --- plus a sixth
regime adding a purely epistemic action channel.

\begin{description}
  \item[R1 --- Passive baseline.] Action stream produced by the analytic
  Bayes-adaptive policy; the model only predicts. Control: H2 predicts no Self
  representation, since the model has no causal role in the action stream. (This
  is \citep{shai2024belief} with an action channel the model does not own.)

  \item[R2 --- Active, single-actor, supervised.] The model produces actions
  autoregressively. Two sub-regimes isolate where the training signal must sit:
  \textbf{(2a)} cross-entropy at both action and observation positions;
  \textbf{(2b)} cross-entropy at observation positions only (pure world model
  conditioned on its own action history). The 2a/2b contrast is the cleanest
  single experiment in the project: it asks whether the causal-sampling
  structure of self-produced tokens \emph{alone} induces a token-type Self
  (2b), or whether explicit gradient at action positions is required (2a).

  \item[R3 --- Active, single-actor, RL.] As R2 but trained on trajectory
  reward rather than next-token CE --- the LLM RL-post-training analogue.

  \item[R4 --- Active, single-actor, budgeted exploration.] The model takes a
  budget of $K$ action/observation steps before being scored on the next
  observation (the chain-of-thought analogue). Sweep $K\ge1$ to probe how the
  exploration budget shapes representations.

  \item[R5 --- Mixed actors.] Each trajectory is generated by an actor sampled
  from a family $\{\pi^{(0)},\dots,\pi^{(M)}\}$ --- the model itself
  ($\pi^{(0)}$), a random policy, the analytic Bayes-adaptive policy, an earlier
  checkpoint, and a deliberately distinct policy. The model is not told which
  actor produced a trajectory and is trained as in 2b. To predict observations
  it must infer the actor from action statistics, motivating an actor-identity
  Self. Holding the World fixed while varying the actor is the cleanest test of
  a genuine causal Self/World separation. This is the toy analogue of the LLM
  story: a predictor over many voices in pretraining that becomes the privileged
  producer of one in post-training.

  \item[R6 --- Epistemic queries.] Add a purely epistemic action channel: the
  model may spend a budget of query tokens, each returning a fixed
  coarse-graining (epimorphic image) $\varphi_i(w_t)$ of the current World state
  \emph{without disturbing it}. Spending the budget well requires consulting the
  quotient the model is most uncertain about. Targets the epistemic self
  (H2c). The behavioural form is deflationary (optimal querying is a readout of
  $b_W$, achievable by a pure world model); the mechanistic form is not. The
  dissociation and intervention tests that break the equivalence are specified
  in \texttt{strand\_three.tex}.
\end{description}

\paragraph{Cross-world (meta-learning) variant.}
\label{sec:metalearn}
Any active regime can be run over a family of SWITCHes sharing
$(\mathcal{W},\mathcal{O})$ but with $T,E$ resampled from a meta-distribution,
with a held-out SWITCH at evaluation. Prediction: the World representation
adapts in-context to the new task while the Self representation stays stable ---
the self as a fixed actor in a volatile world. This variant is also the
precondition for the nested-belief / theory-of-mind geometry of
\texttt{strand\_two.tex} (it requires \emph{stateful} actors in the R5 family).

\section{Measurements}
\label{sec:measure}

Each measurement is a procedure on residual-stream activations $h_\ell^t$ with a
metric and, wherever possible, the analytic comparator from \S\ref{sec:switch}.
Probes are linear unless stated; every probe is reported against a control
(position-only / token-only / shuffled-label baseline) so that ``decodable''
means ``decodable beyond what position and surface token already give''.

\begin{description}
  \item[M1 --- Belief geometry (tests H1).] Fit a linear map $h_\ell^t\mapsto
  \hat b_W^t$ to the analytic posterior; report $\mathrm{KL}(b_W^t\,\|\,\hat
  b_W^t)$ and the recovered MSP geometry (cf.\ \citep{shai2024belief}). Key
  sub-test: does $\hat b_W^t$ implement the action-conditioned update
  \eqref{eq:actorgate}--\eqref{eq:push}, i.e.\ condition correctly on the
  model's own sampled action? A positive result is reflexive self-knowledge of
  the belief.

  \item[M2 --- Factorisation (tests H1$+$H2).] Identify the Self and World probe
  subspaces and test (i) independent decodability of $b_S$ and $b_W$,
  (ii) product structure $b(s,w)\approx b_S(s)\,b_W(w)$ via $\mathrm{KL}$ of the
  decoded joint to the product of decoded marginals, (iii) subspace separation
  via principal angles between the Self and World probe directions. Capacity
  control: vary model width and confirm the factorisation is not forced by a
  bottleneck.

  \item[M3 --- Token-type self (tests H2a).] Decode action-position vs.\
  observation-position role; measure action-position entropy collapse. Predicted
  in R2--R4; sharpest in 2a. The 2a-vs-2b contrast is the decisive sub-test.

  \item[M4 --- Actor-identity self (tests H2b).] Decode actor identity $c$ from
  the residual stream; report AUC. Predicted high in R5, near chance in R2--R4
  (no contrastive signal there). Off-policy test: present an action stream from
  a held-out policy and check whether observation predictions remain
  Bayes-adaptive (Hypothesis~A: abstract action-conditioned update, generalises)
  or degrade (Hypothesis~B: policy-specific shortcut). See
  \texttt{strand\_two.tex} for the circuit-level reading (shared identity-gated
  update vs.\ privileged self pathway).

  \item[M5 --- Epistemic self (tests H2c).] In R6, regress the analytic marginal
  uncertainties (e.g.\ $H(b_A),H(b_B)$ in a factored World) from the residual
  stream and correlate the model's query choice with the decoded-uncertainty
  argmax. \emph{Miscalibration dissociation}: select histories where the
  analytic $b_W$ and the decoded $\hat b_W$ diverge, and classify queries ---
  introspective access tracks $\hat b_W$'s uncertainty, a compiled heuristic
  tracks neither. Detailed in \texttt{strand\_three.tex}.

  \item[M6 --- Training dynamics.] Run M1--M5 on checkpoints. Predicted
  dissociation: World belief sharpens in all regimes; token-type self emerges in
  R2--R4 (sharpest 2a); actor-identity self emerges specifically in R5;
  epistemic self emerges specifically in R6. The 2a-vs-2b emergence contrast
  tells us which pressure is active.
\end{description}

\section{Interventions (causal tests)}
\label{sec:intervene}

Decodability is necessary but not sufficient; we test that the recovered
structure is causally load-bearing by writing along probe-identified directions
and comparing the effect to the analytic prediction.

\begin{description}
  \item[I1 --- Self-steering.] Write a different Self encoding; the action
  distribution should shift as the SWITCH dynamics predict, without disturbing
  $b_W$.
  \item[I2 --- World-steering.] Write a perturbed World encoding; observation
  predictions and downstream actions should shift as the Bayes-adaptive update
  predicts. A clean I1/I2 dissociation is strong evidence the factorisation is
  real, not merely decodable.
  \item[I3 --- Fiber double-dissociation (R5 + stateful actors).] Perturbing the
  self fiber should move both the model's observation predictions and its own
  action distribution; perturbing a non-self fiber $\Delta_W^{(c)}$ should move
  only the predicted behaviour of actor $c$. (See \texttt{strand\_two.tex}.)
  \item[I4 --- Belief-edit $\to$ query (R6).] Edit $\hat b_W$ to a different
  marginal-entropy profile and check the query distribution shifts as the edited
  belief predicts --- a causal pathway from belief representation to epistemic
  action, i.e.\ introspective access. (See \texttt{strand\_three.tex}.)
\end{description}

\section{Predicted outcomes}

\begin{center}
\small
\begin{tabular}{lcccc}
\toprule
Regime & M1 belief & M3 token-type self & M4 actor-identity & M5 epistemic self \\
\midrule
R1 passive            & yes & no            & no  & --- \\
R2a active (both)     & yes & yes (sharpest)& no  & --- \\
R2b active (obs only) & yes & ? (key test)  & no  & --- \\
R3 RL                 & yes & yes           & no  & --- \\
R4 budgeted           & yes & yes           & no  & --- \\
R5 mixed actors       & yes & yes           & yes & --- \\
R6 epistemic queries  & yes & yes           & no  & ? (key test) \\
\bottomrule
\end{tabular}
\end{center}

\noindent The two ``?'' cells are the project's pivotal experiments: whether a
token-type self emerges in 2b without gradient at action positions (M3), and
whether an epistemic self with causal introspective access emerges in R6 (M5,
I4). Either outcome of each is informative.

\section{Theoretical targets}

The two strands below specify the geometric object each probe aims at, so that
M2--M5 and I1--I4 have concrete targets rather than open-ended ``find
structure'' goals.

% =====================================================================
% Strand Two: the geometry of the internal Self/World representations.
% Intended to slot into the "Analysis of Trained Networks" section,
% expanding the "Belief geometry" and "Self-recognition" paragraphs.
% Reuses only citation keys already present in references.bib.
% =====================================================================

\subsection*{The geometry of the Self/World factorisation}

The analyses above ask \emph{whether} the SWITCH factorisation is recovered; here we
make precise \emph{what geometric object} we expect to find, so that the probes and
interventions have a concrete target. In the passive predictor of
\citep{shai2024belief} the object is well understood: a belief is a point in the
probability simplex $\Delta_W$ over the World hidden states, the Bayes-optimal belief
given a history is a single point in $\Delta_W$, and the set of all reachable beliefs
is the (typically fractal) mixed-state presentation (MSP), the attractor of the
iterated function system whose maps are the observation-conditioned Bayes updates. Two
computations sit on top of this geometry: a \emph{linear} emission read-out giving the
next-token distribution, and a \emph{nonlinear} filtering update that moves the belief
point as each observation arrives, both of which \citep{shai2024belief} recover
linearly from the residual stream.

\paragraph{From one simplex to a product.} The SWITCH hidden state is the pair $(s,w)$,
so the belief lives in the joint simplex over $S\times W$. The central hypothesis,
stated geometrically, is that the reachable-belief manifold collapses onto the
\emph{product} submanifold $\Delta_S\times\Delta_W$ --- that the joint belief factors as
$b(s,w)\approx b_S(s)\,b_W(w)$ with the two marginals independently decodable --- rather
than filling an entangled subset of the joint simplex. This is the precise sense in
which Self and World are ``separable'', and it ties our result to the Simplex line of
work on the representation of independent generative processes. Token parity now selects
between two semantically distinct updates rather than one: an \emph{observation} token
triggers a filtering update that contracts $b_W$ toward consistency with what was seen,
while an \emph{action} token triggers a \emph{control} pushforward
$b_W \mapsto P(\cdot\mid\cdot,a)^{\top} b_W$ that advances the World belief through the
action-conditioned dynamics \citep{bengio1995iohmm}. Recovering two distinct update
maps, selected by parity and acting on a product geometry, is the first concrete target
of this strand.

\paragraph{The Self simplex as a persona-selection model.} In Regimes 2--4 the model is
always the actor, so $b_S$ is a point mass: there is no uncertainty about who is acting
and the Self marginal is degenerate, leaving only the thin token-type (action vs.\
observation) distinction. The Self simplex becomes a non-trivial object precisely in the
mixed-actor Regime 5, where each trajectory is generated by one of a family of policies
$\{\pi^{(1)},\dots,\pi^{(M)}\}$, one of which ($\pi^{(0)}$) is the model itself. Here
$\Delta_S$ is the online posterior over \emph{actor identity}: its vertices are the
members of the actor family, and it is updated by the likelihood ratio of each observed
action under the candidate policies --- exactly a filtering process on a ``which-actor''
latent. We therefore expect a second, lower-dimensional mixed-state geometry --- a
persona-selection MSP --- to appear alongside the World MSP, linearly decodable from the
residual stream in Regime 5 but absent in Regimes 2--4.

\paragraph{Nested belief and theory of mind.} The persona simplex does more than label
trajectories: it gates the World update. An action is evidence about the World state
only insofar as the acting policy is World-state-dependent, and how strongly depends on
which actor produced it. The action-token update on $b_W$ therefore carries an
actor-dependent Bayesian factor \emph{before} the pushforward,
\[
  b_W(w)\;\propto\; b_W(w)\,\pi^{(c)}(a\mid w)\quad\text{(actor-gated)},
  \qquad\text{then}\qquad
  b_W \;\mapsto\; P(\cdot\mid\cdot,a)^{\top} b_W,
\]
so that a random-actor action is treated as uninformative (pushforward only) while an
optimal-actor action drives a strong belief update. When the actors are themselves
\emph{stateful} --- e.g.\ Bayes-adaptive agents carrying their own World belief ---
predicting their actions requires the model to simulate \emph{that actor's} belief, a
separate point in a belief simplex that need not coincide with the model's own. The full
latent geometry is then a \emph{bundle}: over each non-self node $c$ of $\Delta_S$ sits a
fiber simplex $\Delta_W^{(c)}$ encoding the model's representation of actor $c$'s
beliefs, distinct both from the ground-truth World state and from the model's own
belief. Representing another agent's belief as separable from one's own is first-order
theory of mind; the SWITCH with stateful actors forces it as a computational requirement
of observation prediction, and it is jointly testable with the World geometry on the
same analytic ground truth.

\paragraph{The self as a degenerate fiber.} The self node is distinguished not by its
location in $\Delta_S$ but by the status of its fiber: over the self node the belief
fiber is \emph{identified} with the model's own World belief $b_W$ rather than separately
inferred, because the self acts on its own belief by construction. This identification is
the geometric content of the reflexive self-recognition reported behaviourally in
post-trained models \citep{betley2025tellme,panickssery2024self} and in our forthcoming
work with Lindsey. It is also what makes the Self component causally load-bearing: a
perturbation written into the self fiber should move both the model's observation
predictions and its own action distribution, whereas a perturbation to a non-self fiber
$\Delta_W^{(c)}$ should move only the predicted behaviour of actor $c$. This double
dissociation, layered on top of the Self-steering interventions already described, is the
sharpest causal signature of a genuine self representation.

\paragraph{Circuit-level predictions and the off-policy test.} Two implementations are
distinguishable and worth separating experimentally. Under a \emph{shared, identity-gated}
update circuit, a single mechanism performs the actor-conditioned belief update with the
$\Delta_S$ vector entering as a gating input; steering $\Delta_S$ should then continuously
deform the observed update rule, from ``ignore the action as evidence'' to ``treat it as
strong evidence''. Under a \emph{privileged self pathway}, the self is handled by the
first-person generation circuit --- read off, not inferred --- rather than by the shared
theory-of-mind module, and only the non-self nodes route through the latter.
Distinguishing these is the mechanistic content of the off-policy generalisation test
already proposed: a shared identity-gated bundle should interpolate to a held-out actor
of a known type (Hypothesis A, generalisation), whereas per-actor shortcuts through the
joint (action, observation) statistics should not (Hypothesis B). Two design notes follow
directly: the nested fiber structure arises only if the actor family contains stateful
policies (memoryless actors require actor identity but no fiber), and the actors should
remain first-order World modellers in the toy --- recursive actor-models-actor structure
is the layered self-simulation we defer to the language-model extension.


% =====================================================================
% Strand Three: epistemic actions and introspective access.
% Regime 6: the model may query lumped quotients (epimorphic images)
% of the World HMM. The analysis separates genuine introspective
% access to the model's own belief state from first-order competence
% that merely uses it.
% =====================================================================

\subsection*{Epistemic actions: queryable quotients and introspective access}

The regimes above vary \emph{who} acts; this strand varies \emph{what acting is
for}. Following the classical distinction of \citep{kirsh1994epistemic}, the
SWITCH actions of Regimes 2--5 are simultaneously \emph{pragmatic} (they
transition the World state) and \emph{epistemic} (the response is informative
about it). Regime 6 adds a channel that is purely epistemic: it yields
information about $w$ without disturbing it.

\paragraph{Setup (Regime 6).} Fix a family of epimorphisms
$\varphi_i : W \to W^{(i)}$ onto smaller HMMs, chosen so that each quotient is a
lumpable coarse-graining of the World chain \citep{kemeny1960finite}. At
designated positions the model may emit a query token $q_i$; the environment
answers with the current coarse state $y = \varphi_i(w_t)$, leaving $w_t$
untouched. Queries are budgeted (a fixed number per trajectory, or a small
reward cost), so the model must choose \emph{which} quotient to consult. The
cleanest instantiation factors the World outright, $W \cong A \times B$, with
the $\varphi_i$ the two projections: the model holding belief $b_W$ must decide
whether it is more ignorant about the $A$ factor or the $B$ factor. The
intuition motivating the regime is that spending the budget well requires the
model to \emph{know what it knows}: a metacognitive judgement about its own
belief state, made before any new evidence arrives.

\paragraph{The deflationary observation.} We state at the outset the objection
that this intuition does not survive unmodified, because the objection shapes
the entire analysis plan. For an ideal Bayes-adaptive agent the expected
information gain of query $i$ is the entropy of the pushforward of its current
belief, $\mathrm{EIG}_i = H(\varphi_{i\,\sharp}\, b_W)$: the answer is a
deterministic function of $w$, so the mutual information $I(w;\varphi_i(w))$
reduces to the entropy of the coarse marginal. The optimal query policy
$\arg\max_i H(\varphi_{i\,\sharp}\, b_W)$ is therefore a functional of the
first-order belief --- one more readout of $b_W$, requiring no representation
\emph{of} the belief over and above the belief itself. Optimal query behaviour
alone is thus compatible with a pure world model. This is precisely the
deflationary critique long aimed at animal and infant metacognition: opt-out
and help-seeking paradigms \citep{hampton2001monkeys,goupil2016infants} admit
first-order reconstructions in which an uncertainty signal is \emph{used}
without being \emph{represented} \citep{carruthers2008skeptical}, and
behavioural evidence alone has not resolved that dispute. The toy setting
cannot evade the equivalence at the behavioural level --- but it can break it
at the mechanistic level, in two ways.

\paragraph{Escape one: dissociation under miscalibration.} The behavioural
equivalence holds only for the ideal agent, for whom ``my uncertainty'' is a
deterministic function of the observable history. A finite trained model's
realised belief $\hat b_W$ --- the object our probes decode --- deviates from
the analytic posterior $b_W$ somewhere: off the training distribution, on long
or ambiguous histories, at capacity limits. There the hypotheses come apart.
\emph{Introspective access} predicts queries that track the model's own
miscalibration, $\arg\max_i H(\varphi_{i\,\sharp}\, \hat b_W)$: the model asks
about what \emph{it} is uncertain of, even where an ideal observer would not
be. A \emph{compiled heuristic} --- a query policy fit to the typical
uncertainty profile of surface history features --- instead approximates the
ideal policy on-distribution and tracks neither posterior where they diverge.
Because both $b_W$ (analytically) and $\hat b_W$ (by the belief-geometry
probes) are available to us, we can locate histories on which they disagree and
classify every query the model makes. Knowing-what-you-know, on this
operationalisation, is not being right about the world; it is being right about
one's own errors.

\paragraph{Escape two: causal introspection.} The sharper test is
interventional and reuses machinery this proposal already builds. Once the
World-belief subspace is identified, we edit $\hat b_W$ in the residual stream
--- writing in a belief with a different marginal-entropy profile --- and ask
whether the query distribution shifts as the edited belief predicts. A positive
result establishes a causal pathway from the belief representation to the query
choice: the metacognitive act \emph{reads} the model's own internal state
rather than recomputing a shortcut from the history. This is the toy-model
analogue of concept-injection tests of introspection in language models, and
connects directly to the on-policy-detection findings in our forthcoming work
with Lindsey. We note the regress this stops: one can always redescribe
``reading one's own belief state'' as more first-order computation, but at that
point the deflation has no further content --- a circuit that consumes the
belief representation \emph{as} a representation, tracks its idiosyncratic
errors, and is causally load-bearing for behaviour is everything introspective
access could be in any physical system.

\paragraph{The epistemic self and its geometry.} If the introspective readout
exists, Regime 6 adds a third notion of self to the two already in play
(token-type and actor-identity): an \emph{epistemic self}, a represented
quantity whose content is a property of the model's own belief state rather
than of the World. In the factored instantiation the predicted object is
concrete: scalar readouts of the marginal entropies $H(b_A)$ and $H(b_B)$,
linearly available upstream of the query decision. These quantities sit on the
Self side of the Markov blanket while being \emph{about} the World --- the
first genuinely second-order representation the SWITCH demands, and the
self-directed counterpart of the other-directed belief fibers of the previous
strand. This symmetry raises a sharp circuit-level question, parallel to the
privileged-self-pathway question raised there: does the self-uncertainty
readout route through the same machinery that (with stateful actors) represents
\emph{other} agents' beliefs, or through a privileged direct pathway? The
former is the mechanistic form of the thesis that metacognition is self-directed
mindreading \citep{carruthers2009mindreading}; the SWITCH with both stateful
actors and queryable quotients makes that long-standing question empirical.

\paragraph{Design notes.} (i) Query timing is curriculum-staged: fixed query
slots first, so the model chooses only \emph{which} quotient to consult, with
learned timing and explicit query costs as a second phase. This keeps the
regime free of the cold-start problem that makes open-ended reasoning training
expensive. (ii) The quotient family should be designed so that the identity of
the most informative quotient varies across histories (the product projections
already guarantee this, as uncertainty migrates between factors); otherwise a
static preference over oracles is optimal and the regime tests nothing.
(iii) The regime reuses the Regime 2b training infrastructure, and its analytic
comparators --- filtered posterior, pushforward entropies, value of information
--- are cheap to compute. (iv) A natural extension, deferred on cost grounds,
replaces the external oracle with free-form internal rollouts in which the
model generates both halves of a simulated exchange (the chain-of-thought
analogue); the quotient-oracle version isolates the metacognitive question
without first having to train reasoning from scratch.


\section{Phasing and open implementation questions}

\paragraph{Phase 0 (de-risking, no new training).} Validate the SWITCH sampler
and the analytic posterior against Monte-Carlo estimates. Reproduce
\citep{shai2024belief} on R1. Run the interventional/counterfactual replay probe
(replay a prefix, swap $a_t\to a'$, decode which posterior the continuation is
conditioned on) on whatever active checkpoint exists first --- a zero-cost
pilot for whether action-conditioned belief is present at all.

\paragraph{Phase 1.} R1, R2a/2b, R5; measurements M1--M4, M6; interventions
I1--I2. This is the deliverable core (H1 and the token-type/actor-identity halves
of H2).

\paragraph{Phase 2.} R3, R4, R6; M5, I3--I4; cross-world meta-learning variant;
stateful actors and the nested-belief geometry. Stretch.

\paragraph{Phase 3 (conditional).} LLM extension to code-execution-augmented
reasoning, where the Python interpreter plays World and the model's code blocks
play Self-emitted actions; paired base/post-trained checkpoints from the
rStar-Math / Qwen-Math-with-code family. Detailed in the grant version.

\paragraph{Open implementation questions.}
\begin{itemize}
  \item \emph{Lumpability for R6:} the quotients $\varphi_i$ should be genuine
  HMM epimorphisms (lumpable coarse-grainings) so the query answer is a clean
  function of $w$; non-lumpable projections make the answer stochastic given $w$
  and weaken the deflation from the start --- a deliberate scope choice to make,
  not a default.
  \item \emph{Self-state observability:} whether to define $g$ (the self-state
  update) so that $b_S$ is non-trivial outside R5/R6, or to accept that the
  single-actor Self is degenerate (a point mass) and only token-type structure
  is available there. The latter is the honest reading and shapes what a null in
  R2--R4 means: ``no represented self'', not ``no self''.
  \item \emph{Reward shaping for R3/R4:} cold-start for RL-discovered behaviour;
  whether a supervised warm-start (distilling the analytic policy) is needed
  before outcome-only training, mirroring the SFT$\to$RL LLM pipeline.
  \item \emph{Probe family:} linear first; if linear probes fail where the
  analytic posterior says structure exists, escalate to low-rank/MLP probes with
  the corresponding control baselines, and treat the gap as itself a finding.
\end{itemize}

\bibliography{references}

\end{document}
